\documentclass[11pt]{article}
\usepackage{preamble}
\renewcommand{\answer}[1]{}

\begin{document}

\ladderheading{AP Statistics \textperiodcentered\ Unit 5 \textperiodcentered\ Topic 5.1 \textperiodcentered\ B: 2027-02-05 \textperiodcentered\ E: 2027-03-31}{One Plot, Two Reporting Strategies}

\focusbanner{TODAY'S PROBLEM}

Suppose a test kitchen measures dough volume for ten batches at assigned times after mixing. The scatterplot shows the paired measurements.\par\smallskip \textbf{Jules:} Use a straight-line summary because volume generally increases with time. Treat the 50-minute batch as ordinary variation.\par \textbf{Ravi:} Decide the form from how the rate changes, and report the 50-minute batch separately if it departs from nearby values.\par Selected measurements (minutes, mL): $(0,420)$, $(20,590)$, $(70,770)$, $(90,785)$.\par
\begin{center}
\begin{tikzpicture}
\begin{axis}[
  width=0.70\linewidth,
  height=1.89in,
  xmin=0, xmax=90, ymin=380, ymax=820,
  xlabel={Minutes after mixing}, ylabel={Dough volume (mL)}, grid=major,
  tick label style={font=\small}, label style={font=\small}
]
\addplot[only marks, mark=*, mark size=1.8pt, color=dayblue] coordinates {
  (0,420)
  (10,510)
  (20,590)
  (30,650)
  (40,695)
  (50,600)
  (60,750)
  (70,770)
  (80,780)
  (90,785)
};
\end{axis}
\end{tikzpicture}
\end{center}
\begin{firsttakebox}{FIRST TAKE --- before the video (5 min)}
Which reporting strategy seems more defensible, Jules's or Ravi's? Cite one graph feature.
\workspace{0.85in}
\end{firsttakebox}

\begin{callout}{\IconBook\ BUILD A COMPLETE SCATTERPLOT DESCRIPTION (keep on the page)}{calloutgreen}
Put the explanatory variable on the $x$-axis and the response on the $y$-axis.
Describe direction, form, strength, and unusual features in context. Linear form has
an approximately constant rate; nonlinear form may curve or change rate. Report an
unusual point without letting that one point erase the pattern followed by the rest.
\end{callout}

\newpage
\textbf{Finish after the video:}\par\smallskip

\dokbadge{1}~\textbf{(a)} Name the explanatory and response variables, with units.\par
\workspace{0.55in}

\dokbadge{2}~\textbf{(b)} Find the volume increase from 0 to 20 minutes and from 70 to 90 minutes. Is the 50-minute batch above or below the nearby batches?\par
\workspace{1.1in}

\dokbadge{3}~\textbf{(c)} Whose reporting strategy is supported? Describe the direction, form, strength, and unusual point in context, using your two increases. Explain how a straight-line summary could mislead someone predicting later growth. Write 3--4 sentences.\par
\workspace{1.7in}

\begin{sentenceframebox}\raggedright\textbf{Frame:} I choose \blank[0.7in]: the pattern is \blank[1.0in] in direction, \blank[1.0in] in form, and \blank[1.0in] in strength.\end{sentenceframebox}

\begin{sentenceframebox}\raggedright\textbf{Frame:} The two increases, \blank[0.6in] and \blank[0.6in] mL, show \blank[1.6in]; a straight line could mislead by \blank[1.8in].\end{sentenceframebox}

\begin{summaryexitbox}{EXIT --- turn this in}
One thing the video changed about my first take: \hrulefill\par
\workspace{0.4in}
\end{summaryexitbox}

\end{document}
